\documentclass{article}
\usepackage{amsmath,amssymb,amsthm}
\usepackage{algorithm,algorithmic}
\usepackage{bm}
\usepackage{enumitem}
\newtheorem{assumption}{Assumption}
\newtheorem{theorem}{Theorem}
\newtheorem{lemma}{Lemma}
\newcommand{\prox}{\operatorname{prox}}
\newcommand{\grad}{\nabla}
\newcommand{\R}{\mathbb{R}}
\newcommand{\EE}{\mathbb{E}}
\newcommand{\norm}[1]{\left\|#1\right\|}
\begin{document}

\section*{Proximal Gradient Descent (PGD) for Non‑Convex Smooth $f$}

\section*{Mathematical Formulation}
Let $f:\R^{d}\to\R$ be continuously differentiable (possibly non‑convex) and let $g:\R^{d}\to\R\cup\{+\infty\}$ be a proper, closed, convex regulariser.  
Define the composite objective
\[
F(x)\;:=\;f(x)+g(x),\qquad x\in\R^{d}.
\]

Given a stepsize $\eta>0$, the proximal gradient iteration reads
\begin{equation}
\boxed{%
x_{k+1}\;=\;\prox_{\eta g}\!\bigl(x_{k}-\eta\,\grad f(x_{k})\bigr)},
\qquad k=0,1,2,\dots
\label{eq:pgd}
\end{equation}
where the proximal operator of $g$ is
\[
\prox_{\eta g}(v)\;:=\;\arg\min_{u\in\R^{d}}
\Bigl\{ g(u)+\frac{1}{2\eta}\norm{u-v}^{2}\Bigr\}.
\]

For convenience we introduce the \emph{proximal gradient mapping}
\[
G_{\eta}(x)\;:=\;\frac{1}{\eta}\Bigl(x-\prox_{\eta g}\!\bigl(x-\eta\grad f(x)\bigr)\Bigr).
\tag{1}
\]

\section*{Pseudocode}
\begin{algorithm}[H]
\caption{Proximal Gradient Descent}
\begin{algorithmic}[1]
\REQUIRE Initial point $x_{0}\in\R^{d}$, stepsize $\eta>0$, max iterations $K$.
\FOR{$k=0,\dots,K-1$}
    \STATE $v_{k}\gets x_{k}-\eta\,\grad f(x_{k})$
    \STATE $x_{k+1}\gets \prox_{\eta g}(v_{k})$
\ENDFOR
\ENSURE $x_{K}$
\end{algorithmic}
\end{algorithm}

\section*{Dry Run Example}
Consider the scalar problem
\[
f(w)=(w-3)^{2},\qquad g\equiv0,\qquad \eta=0.1,\qquad w_{0}=0.
\]
Since $g\equiv0$, $\prox_{\eta g}(v)=v$.
\begin{center}
\begin{tabular}{c|c|c|c}
$k$ & $w_{k}$ & $\grad f(w_{k})=2(w_{k}-3)$ & $w_{k+1}=w_{k}-\eta\grad f(w_{k})$\\\hline
0 & $0$ & $-6$ & $0-0.1(-6)=0.6$\\
1 & $0.6$ & $2(0.6-3)=-4.8$ & $0.6-0.1(-4.8)=1.08$\\
2 & $1.08$ & $2(1.08-3)=-3.84$ & $1.08-0.1(-3.84)=1.464$\\
\end{tabular}
\end{center}
Objective values: $F(w_{0})=9$, $F(w_{1})=(0.6-3)^{2}=5.76$, $F(w_{2})\approx4.33$, $F(w_{3})\approx3.56$.
Hence the method reduces the objective even though $f$ is non‑convex (here actually convex).

\newpage

\section*{Assumptions}
\begin{assumption}[Smoothness of $f$]\label{as:smooth}
$f$ has $L$‑Lipschitz continuous gradient, i.e.
\[
\forall x,y\in\R^{d}:\quad
\norm{\grad f(x)-\grad f(y)}\le L\norm{x-y}.
\tag{2}
\]
\end{assumption}
\begin{assumption}[Convexity of $g$]\label{as:convex}
$g$ is proper, closed and convex; consequently the proximal operator $\prox_{\eta g}$ is \emph{firmly non‑expansive}:
\[
\norm{\prox_{\eta g}(u)-\prox_{\eta g}(v)}^{2}
\le \langle u-v,\prox_{\eta g}(u)-\prox_{\eta g}(v)\rangle,
\qquad\forall u,v\in\R^{d}.
\tag{3}
\]
\end{assumption}
\begin{assumption}[Bounded below]\label{as:bounded}
The composite objective $F$ is bounded below:
\[
F^{\inf}:=\inf_{x\in\R^{d}}F(x)>-\infty .
\tag{4}
\]
\end{assumption}

\section*{Main Convergence Theorem}
\begin{theorem}\label{thm:main}
Assume \ref{as:smooth}–\ref{as:bounded} and choose a stepsize
\[
0<\eta\le\frac{1}{L}.
\tag{5}
\]
Then the iterates generated by \eqref{eq:pgd} satisfy for every $K\ge1$
\[
\frac{1}{K}\sum_{k=0}^{K-1}\norm{G_{\eta}(x_{k})}^{2}
\;\le\;
\frac{2}{\eta K}\bigl(F(x_{0})-F^{\inf}\bigr).
\tag{6}
\]
Consequently
\[
\min_{0\le k\le K-1}\norm{G_{\eta}(x_{k})}^{2}
\;\le\;
\frac{2}{\eta K}\bigl(F(x_{0})-F^{\inf}\bigr)
=O\!\left(\frac{1}{K}\right).
\]
Thus any limit point of $\{x_{k}\}$ is a \emph{stationary point} of $F$, i.e.
$0\in\grad f(\bar x)+\partial g(\bar x)$.
\end{theorem}

\section*{Theoretical Properties}
\begin{itemize}[leftmargin=*]
\item \textbf{Stability.} Condition \eqref{5} guarantees that the descent inequality \eqref{eq:descent} holds; larger $\eta$ may violate it and cause divergence.
\item \textbf{Hyper‑parameter Sensitivity.} The bound \eqref{6} scales as $1/\eta$; a smaller stepsize yields a tighter guarantee on $\norm{G_{\eta}}$ but slows practical progress.
\item \textbf{Comparison to Gradient Descent.} If $g\equiv0$, $G_{\eta}(x)=\grad f(x)$ and \eqref{6} reduces to the classical $O(1/K)$ stationarity rate for gradient descent on smooth non‑convex functions.
\item \textbf{Special Cases.}
    \begin{enumerate}
        \item If $g$ is the indicator of a convex set $\mathcal C$, then $\prox_{\eta g}$ is the Euclidean projection onto $\mathcal C$, and $G_{\eta}$ becomes the projected gradient mapping.
        \item If $g$ is $\lambda\|x\|_{1}$, $\prox_{\eta g}$ is the soft‑thresholding operator; the theorem thus covers the proximal‑gradient method for $\ell_{1}$‑regularised non‑convex loss.
    \end{enumerate}
\end{itemize}

\section*{Preliminaries (Definitions and Lemmas)}
\begin{lemma}[Descent Lemma for $L$‑smooth $f$]\label{lem:descent}
For any $x,y\in\R^{d}$,
\[
f(y)\le f(x)+\langle\grad f(x),y-x\rangle+\frac{L}{2}\norm{y-x}^{2}.
\tag{7}
\]
\end{lemma}
\begin{proof}
Standard consequence of $L$‑Lipschitz gradient; see e.g. \cite{Nesterov2004}. \hfill $\square$
\end{proof}

\begin{lemma}[Firm non‑expansiveness of $\prox$]\label{lem:firm}
If $g$ is convex, then for any $u,v\in\R^{d}$,
\[
\|\prox_{\eta g}(u)-\prox_{\eta g}(v)\|^{2}
\le \langle u-v,\,\prox_{\eta g}(u)-\prox_{\eta g}(v)\rangle.
\tag{8}
\]
\end{lemma}
\begin{proof}
Directly from the definition of the proximal operator as the minimiser of a strongly convex subproblem. \hfill $\square$
\end{proof}

\section*{Main Proof (Step‑by‑Step Derivation)}
We prove Theorem~\ref{thm:main}.  Throughout the proof we write $x_{k+1}$ as in \eqref{eq:pgd} and denote
\[
u_{k}:=x_{k}-\eta\grad f(x_{k}),\qquad
v_{k}:=\prox_{\eta g}(u_{k})=x_{k+1}.
\]

\subsection*{Step 1 – Apply the descent lemma to $f$}
From Lemma~\ref{lem:descent} with $x=x_{k}$ and $y=v_{k}$,
\[
f(v_{k})\le f(x_{k})+\langle\grad f(x_{k}),v_{k}-x_{k}\rangle
+\frac{L}{2}\norm{v_{k}-x_{k}}^{2}.
\tag{9}
\]

\subsection*{Step 2 – Relate $v_{k}-x_{k}$ to the proximal gradient mapping}
By definition of $G_{\eta}$,
\[
v_{k}=x_{k}-\eta G_{\eta}(x_{k})\quad\Longrightarrow\quad
v_{k}-x_{k}=-\eta G_{\eta}(x_{k}).
\tag{10}
\]
Insert \eqref{10} into \eqref{9}:
\[
f(v_{k})\le f(x_{k})-\eta\langle\grad f(x_{k}),G_{\eta}(x_{k})\rangle
+\frac{L\eta^{2}}{2}\norm{G_{\eta}(x_{k})}^{2}.
\tag{11}
\]

\subsection*{Step 3 – Use optimality of the proximal step}
The definition of $v_{k}$ yields the first‑order optimality condition
\[
0\in \partial g(v_{k})+\frac{1}{\eta}(v_{k}-u_{k})
\;\Longrightarrow\;
\frac{1}{\eta}(u_{k}-v_{k})\in\partial g(v_{k}).
\tag{12}
\]
Since $u_{k}=x_{k}-\eta\grad f(x_{k})$, we have
\[
\frac{1}{\eta}(u_{k}-v_{k})=\grad f(x_{k})-G_{\eta}(x_{k})\in\partial g(v_{k}).
\tag{13}
\]

\subsection*{Step 4 – Combine $f$ and $g$}
Add $g(v_{k})$ to both sides of \eqref{11} and use convexity of $g$ together with \eqref{13}:
\[
\begin{aligned}
F(v_{k}) & = f(v_{k})+g(v_{k})\\
&\le f(x_{k})+g(v_{k})
-\eta\langle\grad f(x_{k}),G_{\eta}(x_{k})\rangle
+\frac{L\eta^{2}}{2}\norm{G_{\eta}(x_{k})}^{2}.
\end{aligned}
\tag{14}
\]
Convexity of $g$ gives
\[
g(v_{k})\le g(x_{k})+\langle s_{k},\,v_{k}-x_{k}\rangle,
\quad\text{for any } s_{k}\in\partial g(x_{k}).
\tag{15}
\]
Choose $s_{k}= \grad f(x_{k})-G_{\eta}(x_{k})$ from \eqref{13} (which belongs to $\partial g(v_{k})$, but convexity also yields the same inequality with $v_{k}$ and $x_{k}$ exchanged). Substituting \eqref{15} into \eqref{14} and using $v_{k}-x_{k}=-\eta G_{\eta}(x_{k})$ yields
\[
\begin{aligned}
F(v_{k})
&\le f(x_{k})+g(x_{k})
-\eta\langle\grad f(x_{k}),G_{\eta}(x_{k})\rangle
+\frac{L\eta^{2}}{2}\norm{G_{\eta}(x_{k})}^{2}
\\
&\qquad
+\langle \grad f(x_{k})-G_{\eta}(x_{k}),\, -\eta G_{\eta}(x_{k})\rangle\\[1ex]
&=F(x_{k})-\eta\norm{G_{\eta}(x_{k})}^{2}
+\frac{L\eta^{2}}{2}\norm{G_{\eta}(x_{k})}^{2}.
\end{aligned}
\tag{16}
\]

\subsection*{Step 5 – Descent inequality}
Rearrange \eqref{16}:
\[
F(x_{k})-F(x_{k+1})
\;\ge\;
\Bigl(\eta-\frac{L\eta^{2}}{2}\Bigr)\norm{G_{\eta}(x_{k})}^{2}.
\tag{17}
\]
Because of the stepsize restriction \eqref{5},
\[
\eta-\frac{L\eta^{2}}{2}\;\ge\;\frac{\eta}{2}.
\tag{18}
\]
Insert \eqref{18} into \eqref{17} to obtain the key descent bound
\[
\boxed{%
F(x_{k})-F(x_{k+1})\;\ge\;\frac{\eta}{2}\norm{G_{\eta}(x_{k})}^{2}
}\qquad\text{[Step 1]}.
\tag{19}
\]

\subsection*{Step 6 – Telescoping sum}
Sum \eqref{19} over $k=0,\dots,K-1$:
\[
\sum_{k=0}^{K-1}\norm{G_{\eta}(x_{k})}^{2}
\;\le\;\frac{2}{\eta}\sum_{k=0}^{K-1}\bigl(F(x_{k})-F(x_{k+1})\bigr)
=\frac{2}{\eta}\bigl(F(x_{0})-F(x_{K})\bigr).
\tag{20}
\]
Using the lower bound \eqref{4},
\[
F(x_{K})\ge F^{\inf}\quad\Longrightarrow\quad
\sum_{k=0}^{K-1}\norm{G_{\eta}(x_{k})}^{2}
\le\frac{2}{\eta}\bigl(F(x_{0})-F^{\inf}\bigr).
\tag{21}
\]

\subsection*{Step 7 – Average and minimum stationarity measure}
Dividing \eqref{21} by $K$ yields the average bound \eqref{6}:
\[
\frac{1}{K}\sum_{k=0}^{K-1}\norm{G_{\eta}(x_{k})}^{2}
\le\frac{2}{\eta K}\bigl(F(x_{0})-F^{\inf}\bigr).
\tag{22}
\]
Since the minimum of a set of non‑negative numbers does not exceed the average,
\[
\min_{0\le k\le K-1}\norm{G_{\eta}(x_{k})}^{2}
\le\frac{2}{\eta K}\bigl(F(x_{0})-F^{\inf}\bigr).
\tag{23}
\]

\subsection*{Step 8 – Convergence to a stationary point}
Because the right‑hand side of \eqref{23} tends to $0$ as $K\to\infty$, the sequence $\{G_{\eta}(x_{k})\}$ converges to $0$ in the sense that
\[
\liminf_{k\to\infty}\norm{G_{\eta}(x_{k})}=0.
\]
By definition \eqref{1}, $G_{\eta}(x_{k})\to0$ implies
\[
0\in\grad f(\bar x)+\partial g(\bar x)
\]
for any accumulation point $\bar x$ of $\{x_{k}\}$, i.e. $\bar x$ is a stationary point of $F$.
\hfill $\square$

\section*{Rate Derivation (With Constants)}
From \eqref{23} we explicitly have
\[
\min_{0\le k\le K-1}\norm{G_{\eta}(x_{k})}^{2}
\le \underbrace{\frac{2\bigl(F(x_{0})-F^{\inf}\bigr)}{\eta}}_{C}
\cdot\frac{1}{K},
\]
so the constant $C$ depends linearly on the initial optimality gap and inversely on the stepsize $\eta$.

\section*{Special Cases}
\begin{enumerate}
\item \textbf{Projected Gradient Descent.} If $g=\iota_{\mathcal C}$ (indicator of a convex set), $\prox_{\eta g}$ is the Euclidean projection $P_{\mathcal C}$. The mapping $G_{\eta}$ reduces to the projected gradient mapping, and Theorem~\ref{thm:main} recovers the standard $O(1/K)$ stationarity rate for projected gradient on non‑convex $f$.
\item \textbf{Lasso‑type Regularisation.} For $g(x)=\lambda\|x\|_{1}$, $\prox_{\eta g}$ is soft‑thresholding. The same $O(1/K)$ bound holds, providing a guarantee for sparse non‑convex learning problems.
\item \textbf{Exact Gradient Descent.} When $g\equiv0$, $G_{\eta}(x)=\grad f(x)$ and \eqref{19} simplifies to $f(x_{k})-f(x_{k+1})\ge\frac{\eta}{2}\|\grad f(x_{k})\|^{2}$, the classical result.
\end{enumerate}

\begin{thebibliography}{9}
\bibitem{Nesterov2004}
Y.~Nesterov, \emph{Introductory Lectures on Convex Optimization: A Basic Course}, Kluwer Academic Publishers, 2004.
\end{thebibliography}

\end{document}